%───────────────────────────────────────────────────────────────────────────────
\chapter{Galería de referencia implementada}
\label{ch:galeria}
%───────────────────────────────────────────────────────────────────────────────

% INSTRUCCIÓN PARA LA IA QUE LEE ESTE CAPÍTULO:
% Este capítulo es tu referencia de implementación. Cada sección muestra
% un elemento del template implementado correctamente con coloreado semántico,
% theming correcto y citas bibliográficas reales. Úsalo como fuente de
% copy-paste y como verificación antes de generar cualquier contenido.
%
% REGLAS ABSOLUTAS que este capítulo ejemplifica:
%   1. Nunca Ctp* fuera de catppuccin-palette.tex — todos los colores
%      usan aliases semánticos (env-*, tikz-*, plot-*, eq-*, etc.)
%   2. Las macros \eqop, \equ, \eqk, etc. SIEMPRE dentro de modo matemático
%   3. Cada capítulo termina con \printbibliography[heading=subbibliography]
%   4. Los flotantes (figure, acadtable) NO van dentro de tcolorbox —
%      usar boxfig y boxtab en su lugar

Este capítulo documenta cada elemento del template con ejemplos
funcionales, coloreado correcto y referencias bibliográficas reales.
Está diseñado para ser usado directamente como referencia de implementación.

%───────────────────────────────────────────────────────────────────────────────
\section{Entornos teoremáticos}
\label{sec:gal:entornos}
%───────────────────────────────────────────────────────────────────────────────

% Los ocho entornos tienen firma uniforme:
%   \begin{nombre}{Título de esta instancia}
%   \label{tipo:etiqueta}
%   Contenido con ecuaciones coloreadas...
%   \end{nombre}
%
% NUNCA uses un entorno en lugar de otro por estética.
% Cada uno tiene un rol epistémico preciso — ver README.md.

\subsection{Nivel 1 — Fundamentos}

\begin{definicion}{Espacio de Hilbert}
\label{def:gal:hilbert}
Un \textbf{espacio de Hilbert}~\cite{reed1980}
$(\eqdom{\mathcal{H}}, \langle\cdot,\cdot\rangle)$ es un espacio vectorial
sobre $\mathbb{K}$ con producto interior
$\langle \cdot, \cdot \rangle : \eqdom{\mathcal{H}} \times \eqdom{\mathcal{H}}
\to \mathbb{K}$ tal que la norma inducida
$\|\equ{x}\| = \sqrt{\langle \equ{x}, \equ{x} \rangle}$
hace a $\eqdom{\mathcal{H}}$ completo.
\end{definicion}

\begin{teorema}{Teorema de representación de Riesz~\cite{reed1980}}
\label{thm:gal:riesz}
Sea $\eqdom{\mathcal{H}}$ un espacio de Hilbert y
$\eqfn{\varphi} : \eqdom{\mathcal{H}} \to \mathbb{K}$ un funcional lineal
continuo. Existe un único $\eqk{y} \in \eqdom{\mathcal{H}}$ tal que
\[
  \eqfn{\varphi}(\equ{x}) = \langle \equ{x}, \eqk{y} \rangle
  \qquad \forall\, \equ{x} \in \eqdom{\mathcal{H}}.
\]
Además $\|\eqfn{\varphi}\| = \|\eqk{y}\|$.
\end{teorema}

\begin{proposicion}{Identidad de polarización}
\label{prop:gal:polarizacion}
En un espacio de Hilbert real~\cite{munkres2000}, el producto interior
se recupera de la norma:
\[
  \langle \eqk{x}, \eqk{y} \rangle
  = \frac{1}{4}\left(\|\eqk{x}+\eqk{y}\|^2 - \|\eqk{x}-\eqk{y}\|^2\right).
\]
\end{proposicion}

\subsection{Nivel 2 — Apoyo}

\begin{ejemplo}{$\ell^2$ como espacio de Hilbert}
\label{ex:gal:l2}
El espacio $\eqdom{\ell^2} = \{(\equ{x}_n) \subset \mathbb{R} \mid
\eqop{\sum}_n \equ{x}_n^2 < \infty\}$ con producto interior
$\langle \equ{x}, \equ{y} \rangle = \eqop{\sum}_{n=1}^\infty \equ{x}_n \equ{y}_n$
es un espacio de Hilbert separable~\cite{reed1980}. La base ortonormal
canónica $\{\eqk{e}_n\}$ es una base de Schauder pero no algebraica (Hamel).
\end{ejemplo}

\begin{observacion}{Completitud es indispensable}
El espacio $C[0,1]$ con la norma $\eqdom{L^2}$ tiene producto interior pero
\textit{no} es de Hilbert — su completación es $\eqdom{L^2[0,1]}$.
La convergencia débil $\eqk{e}_n \rightharpoonup 0$ en $\eqdom{\ell^2}$
no implica convergencia en norma: $\|\eqk{e}_n\| = 1$ para todo $n$.
Este fenómeno~\cite{reed1980} es fuente frecuente de errores en argumentos
de compacidad en espacios de dimensión infinita.
\end{observacion}

\begin{fisicaconexion}{Estados cuánticos y espacios de Hilbert~\cite{griffiths2018}}
Los estados de un sistema cuántico son vectores en
$\eqdom{\mathcal{H}} = \eqdom{L^2(\mathbb{R}^3)}$.
La norma $\|\equ{\psi}\|^2 = \eqop{\int}_{\eqdom{\mathbb{R}^3}}
|\equ{\psi}(\eqvec{\mathbf{r}})|^2 \eqdm{d^3r} = 1$
es la condición de normalización. La evolución temporal
$\equ{\psi}(t) = e^{-\eqcon{i}\eqop{\hat{H}}t/\eqcon{\hbar}}\eqk{\psi_0}$
es un operador unitario sobre $\eqdom{\mathcal{H}}$.
\end{fisicaconexion}

\subsection{Nivel 3 — Especiales}

\begin{estrella}{Teorema espectral para operadores autoadjuntos}
\label{star:gal:espectral}
Sea $\eqop{A} : \eqdom{\mathcal{H}} \to \eqdom{\mathcal{H}}$ autoadjunto
acotado~\cite{reed1980}. Existe una medida espectral única
$\eqfn{E} : \mathcal{B}(\mathbb{R}) \to \mathcal{L}(\eqdom{\mathcal{H}})$
tal que
\[
  \eqop{A} = \eqop{\int}_{\sigma(\eqop{A})} \equ{\lambda}\,
  \eqdm{dE}(\equ{\lambda}),
\]
y $\eqfn{f}(\eqop{A}) = \eqop{\int} \eqfn{f}(\equ{\lambda})\,
\eqdm{dE}(\equ{\lambda})$ define un cálculo funcional consistente.
\end{estrella}

\begin{ejerciciocomp}{Verificación numérica del método de potencias}
\label{comp:gal:potencias}
Implementa la iteración de potencias para hallar el eigenvalor dominante
$\equ{\lambda_{\max}}$ de una matriz $\eqten{A} \in \eqdom{\mathbb{R}^{n\times n}}$
SPD aleatoria. Verifica que $|\equ{\lambda}_k - \equ{\lambda_{\max}}|$
converge con tasa $|\eqk{\lambda_2}/\eqk{\lambda_1}|^k$
comparando con los eigenvalores exactos de \texttt{np.linalg.eigh}.
Usa $\eqk{n} \in \{50, 200, 500\}$ y reporta la convergencia en una tabla.
\end{ejerciciocomp}

%───────────────────────────────────────────────────────────────────────────────
\section{Matemáticas semánticas}
\label{sec:gal:ecuaciones}
%───────────────────────────────────────────────────────────────────────────────

% REGLA CRÍTICA PARA LA IA:
% Las macros \eqop, \equ, \eqk, \eqfn, \eqdm, \eqdom, \eqvec, \eqten, \eqcon
% SOLO funcionan dentro de modo matemático: $...$, \[...\], equation, align.
% Fuera de modo matemático el color no se aplica — error frecuente de IA.

\dificultad{4}{5}

\subsection{Los tres ejes combinados: Schrödinger}

% Eje 1 (sintáctico): \eqop para operadores, \eqfn para funciones, \eqdm para medidas
% Eje 2 (epistémico): \equ para incógnitas, \eqk para datos, \eqcon para constantes
% Eje 3 (tipo): \eqvec para vectores, \eqten para tensores
\begin{equation}
\label{eq:gal:schrodinger}
\eqop{i}\eqcon{\hbar}\,\eqop{\frac{\partial}{\partial t}}\equ{\Psi}
= \eqop{\hat{H}}\equ{\Psi}
= \left(\eqop{-}\frac{\eqcon{\hbar}^2}{2\eqk{m}}\eqop{\nabla}^2
  + \eqfn{V}(\eqvec{\mathbf{r}},t)\right)\equ{\Psi}.
\end{equation}

El hamiltoniano $\eqop{\hat{H}}$ y $\eqop{\nabla}^2$ son operadores~\cite{griffiths2018}.
$\equ{\Psi}$ es la incógnita. $\eqfn{V}$ es el potencial (función conocida del
sistema). $\eqcon{\hbar}$ y $\eqk{m}$ son constante y parámetro respectivamente.

\subsection{Dualidad epistémica: física vs. ML}

% Esta es la demostración canónica del eje epistémico.
% La estructura algebraica es idéntica; solo el color cambia.
\textbf{Lectura física}~\cite{griffiths2018} —
parámetros $\eqk{\theta}$ conocidos, estado $\equ{x}$ buscado:
\begin{equation}
\label{eq:gal:physics-reading}
\equ{x^*} = \eqop{\argmin}_{\equ{x}}\;
\eqfn{\mathcal{L}}\!\left(\eqfn{f}(\equ{x};\,\eqk{\theta}),\; \eqk{y}\right).
\end{equation}

\textbf{Lectura ML}~\cite{goodfellow2016} —
dato $\eqk{x}$ conocido, parámetros $\equ{\theta}$ aprendidos:
\begin{equation}
\label{eq:gal:ml-reading}
\equ{\theta^*} = \eqop{\argmin}_{\equ{\theta}}\;
\eqfn{\mathcal{L}}\!\left(\eqfn{f}(\eqk{x};\,\equ{\theta}),\; \eqk{y}\right).
\end{equation}

\subsection{Ecuaciones de Maxwell (eje de tipo)}

\begin{align}
\eqop{\nabla \cdot} \eqvec{\mathbf{E}}
  &= \frac{\eqk{\rho}}{\eqcon{\varepsilon_0}},
  \label{eq:gal:gauss-e}\\
\eqop{\nabla \times} \eqvec{\mathbf{E}}
  &= \eqop{-}\eqop{\frac{\partial}{\partial t}}\eqvec{\mathbf{B}},
  \label{eq:gal:faraday}\\
\eqop{\nabla \cdot} \eqvec{\mathbf{B}}
  &= 0,
  \label{eq:gal:gauss-b}\\
\eqop{\nabla \times} \eqvec{\mathbf{B}}
  &= \eqcon{\mu_0}\!\left(\eqvec{\mathbf{J}}
    + \eqcon{\varepsilon_0}\eqop{\frac{\partial}{\partial t}}\eqvec{\mathbf{E}}\right).
  \label{eq:gal:ampere}
\end{align}

Los campos $\eqvec{\mathbf{E}}$, $\eqvec{\mathbf{B}}$, $\eqvec{\mathbf{J}}$
usan \texttt{\textbackslash eqvec}~\cite{jackson1999}. Las constantes
$\eqcon{\varepsilon_0}$, $\eqcon{\mu_0}$ usan \texttt{\textbackslash eqcon}.
La fuente $\eqk{\rho}$ es un dato del problema (\texttt{\textbackslash eqk}).

\subsection{Macros compuestos}

Proceso de Ornstein-Uhlenbeck~\cite{oksendal2003}:
\begin{equation}
\label{eq:gal:ou}
\eqSDE{X_t}{-\eqk{\theta}\equ{X_t}}{\eqk{\sigma}}.
\end{equation}

Fokker-Planck para la densidad $\equ{p}(\eqvec{\mathbf{x}},t)$~\cite{oksendal2003}:
\begin{equation}
\label{eq:gal:fokker-planck}
\eqpdv{p}{t}
= \eqop{-}\eqop{\nabla_x \cdot}\!\left(\eqfn{f}\,\equ{p}\right)
+ \frac{\eqk{\sigma}^2}{2}\,\eqop{\nabla_x^2}\,\equ{p}.
\end{equation}

Score matching~\cite{hyvarinen2005} (ML — $\equ{\theta}$ es la incógnita):
\begin{equation}
\label{eq:gal:score}
\eqfn{\mathcal{L}}(\equ{\theta})
= \eqcon{\tfrac{1}{2}}\,
  \eqop{\mathbb{E}}_{\eqk{x} \sim \eqdom{p_{\text{data}}}}
  \left[
    \norm{\eqfn{s}_{\equ{\theta}}(\eqk{x})
    - \eqop{\nabla}_{\eqk{x}} \log \eqk{p}(\eqk{x})}^2
  \right].
\end{equation}

Tensor de Einstein ($\eqten{g_{\mu\nu}}$ es el tensor métrico):
\begin{equation}
\label{eq:gal:einstein}
\equ{R_{\mu\nu}} - \eqcon{\tfrac{1}{2}}\eqten{g_{\mu\nu}}\equ{R}
+ \eqcon{\Lambda}\eqten{g_{\mu\nu}}
= \eqcon{\frac{8\pi G}{c^4}}\eqk{T_{\mu\nu}}.
\end{equation}

\subsection{Toggle B/N — versión impresión}

% \eqcolorsfalse desactiva todo el coloreado
% Las macros se vuelven transparentes — no hay que editar ninguna ecuación
{%
\eqcolorsfalse
\begin{equation}
\eqop{i}\eqcon{\hbar}\,\eqop{\frac{\partial}{\partial t}}\equ{\Psi}
= \left(\eqop{-}\frac{\eqcon{\hbar}^2}{2\eqk{m}}\eqop{\nabla}^2
  + \eqfn{V}\right)\equ{\Psi}.
\tag{\ref{eq:gal:schrodinger}$'$}
\end{equation}
}

%───────────────────────────────────────────────────────────────────────────────
\section{Algoritmos con coloreado semántico}
\label{sec:gal:algoritmos}
%───────────────────────────────────────────────────────────────────────────────

% Los algoritmos NO están envueltos en tcolorbox — eso rompe el paginado.
% Las expresiones matemáticas dentro del algoritmo usan coloreado semántico normal.

\begin{algorithm}[H]
  \DontPrintSemicolon
  \SetKwInOut{Input}{Entrada}\SetKwInOut{Output}{Salida}
  \Input{$\eqfn{f}:\eqdom{\mathbb{R}^n}\to\mathbb{R}$,
         $\eqop{\nabla}\eqfn{f}$,
         $\eqk{x_0}$, $\eqk{\eta}$, $\eqk{\varepsilon}$}
  \Output{Mínimo aproximado $\equ{x^*}$}
  $\equ{x} \gets \eqk{x_0}$\;
  \While{$\|\eqop{\nabla}\eqfn{f}(\equ{x})\| > \eqk{\varepsilon}$}{
    $\equ{x} \gets \equ{x} - \eqk{\eta}\,\eqop{\nabla}\eqfn{f}(\equ{x})$\;
  }
  \Return{$\equ{x}$}
  \caption{Descenso de gradiente~\cite{goodfellow2016}}
  \label{alg:gal:gradient}
\end{algorithm}

\begin{algorithm}[H]
  \DontPrintSemicolon
  \SetKwInOut{Input}{Entrada}\SetKwInOut{Output}{Salida}
  \Input{SDE $\eqdm{d}\equ{X} = \eqfn{\mu}(\equ{X},t)\eqdm{dt}
         + \eqfn{\sigma}(\equ{X},t)\eqdm{dW}$,
         $\eqk{X_0}$, $\eqk{\Delta t}$, $\eqk{T}$}
  \Output{Trayectoria $\{\equ{X}_k\}_{k=0}^{\eqk{N}}$}
  $\equ{X} \gets \eqk{X_0}$\;
  \For{$k = 0$ \KwTo $\eqk{N}-1$}{
    $t \gets k\,\eqk{\Delta t}$\;
    $\equ{\xi} \gets \mathcal{N}(0,1)$\tcp*{ruido gaussiano}
    $\equ{X} \gets \equ{X}
      + \eqfn{\mu}(\equ{X},t)\,\eqk{\Delta t}
      + \eqfn{\sigma}(\equ{X},t)\sqrt{\eqk{\Delta t}}\,\equ{\xi}$\;
  }
  \Return{trayectoria}
  \caption{Euler-Maruyama~\cite{oksendal2003}}
  \label{alg:gal:em}
\end{algorithm}

%───────────────────────────────────────────────────────────────────────────────
\section{Código fuente}
\label{sec:gal:codigo}
%───────────────────────────────────────────────────────────────────────────────

% Siempre especificar style=lenguaje.
% Para código LaTeX del template, usar style=latex.

\begin{lstlisting}[style=python,
  caption={Descenso de gradiente estocástico con momentum~\cite{goodfellow2016}},
  label={lst:gal:sgd}]
import numpy as np

def sgd_momentum(grad_fn, theta0, lr=1e-3, beta=0.9, n_iter=1000):
    """SGD con momentum de Polyak."""
    theta = theta0.copy()
    v     = np.zeros_like(theta)
    for _ in range(n_iter):
        g  = grad_fn(theta)
        v  = beta * v + (1 - beta) * g
        theta -= lr * v
    return theta
\end{lstlisting}

\begin{lstlisting}[style=julia,
  caption={Iteración de potencias en Julia~\cite{leveque2007}},
  label={lst:gal:power}]
function power_iteration(A::AbstractMatrix; tol=1e-10, maxiter=1000)
    n      = size(A, 1)
    v      = normalize(randn(n))
    lambda = zero(eltype(A))
    for _ in 1:maxiter
        w      = A * v
        lambda = dot(v, w)         # cociente de Rayleigh
        v      = normalize(w)
    end
    return lambda, v
end
\end{lstlisting}

\begin{lstlisting}[style=latex,
  caption={Patrón correcto de ecuación semántica},
  label={lst:gal:eq-pattern}]
% Siempre dentro de modo matematico
\[
  \eqop{i}\eqcon{\hbar}\,\eqop{\partial_t}\equ{\Psi}
  = \eqop{\hat{H}}\equ{\Psi}
\]

% En tabla (celda con modo matematico)
Operador  & $\eqop{\nabla}$    & Azul     \\
Incognita & $\equ{\psi}$       & Peach    \\
Dato      & $\eqk{m}$          & Mauve    \\

% INCORRECTO: fuera de modo matematico - no colorea
\eqop{\nabla}
\end{lstlisting}

%───────────────────────────────────────────────────────────────────────────────
\section{Tablas}
\label{sec:gal:tablas}
%───────────────────────────────────────────────────────────────────────────────

% acadtable: flotante en flujo normal de texto → \textwidth
% boxtab: no-flotante dentro de tcolorbox → \linewidth
% Columnas S de siunitx: encabezados SIEMPRE en \multicolumn{1}{c}{\thead{...}}
% Celdas con texto en col S: \multicolumn{1}{c}{texto}

\begin{acadtable}[H]{tab:gal:optimizadores}{Métodos de optimización de primer orden~\cite{goodfellow2016}}
  \begin{tabularx}{\textwidth}{l Y Y l}
    \toprule
    \headerrow
    \thead{Método} & \thead{Convergencia} & \thead{Memoria} & \thead{Hessiano} \\
    \midrule
    Gradiente      & $\mathcal{O}(1/k)$          & $\mathcal{O}(n)$   & No \\
    \rowcolor{table-row-alt}
    Momentum       & $\mathcal{O}(1/k^2)$         & $\mathcal{O}(n)$   & No \\
    L-BFGS         & Superlineal                   & $\mathcal{O}(mn)$  & Aprox. \\
    \rowcolor{table-row-alt}
    Adam           & $\mathcal{O}(1/\sqrt{k})$    & $\mathcal{O}(n)$   & No \\
    Newton         & Cuadrática                    & $\mathcal{O}(n^2)$ & Exacto \\
    \bottomrule
  \end{tabularx}
\end{acadtable}

\begin{acadtable}[H]{tab:gal:convergencia}{Convergencia de RK4 para $\equ{y}' = \eqop{-}\eqk{\lambda}\equ{y}$~\cite{leveque2007}}
  \begin{tabularx}{\textwidth}{Y S[table-format=1.2e-1] S[table-format=1.4e-2] S[table-format=1.2]}
    \toprule
    \headerrow
    \multicolumn{1}{Y}{\thead{$\eqk{h}$}} &
    \multicolumn{1}{c}{\thead{$\|\equ{e}\|_\infty$}} &
    \multicolumn{1}{c}{\thead{Razón}} &
    \multicolumn{1}{c}{\thead{Orden emp.}} \\
    \midrule
    $1/4$  & 2.31e-3 & \multicolumn{1}{c}{---} & \multicolumn{1}{c}{---} \\
    $1/8$  & 1.48e-4 & 0.0641 & 3.96 \\
    $1/16$ & 9.33e-6 & 0.0630 & 3.99 \\
    $1/32$ & 5.84e-7 & 0.0626 & 4.00 \\
    \bottomrule
  \end{tabularx}
  \tablenote{*}{$\eqk{\lambda}=1$, $\equ{y}(0)=1$, $t\in[0,1]$.}
\end{acadtable}

% boxtab dentro de un entorno tcolorbox:
\begin{definicion}{Media y varianza muestral}
\label{def:gal:estadisticos}
Dado un conjunto $\{\eqk{x_i}\}_{i=1}^{\eqk{n}}$ de observaciones i.i.d.:
\begin{boxtab}{tab:gal:stats}{Estimadores muestrales}
  \begin{tabularx}{\linewidth}{l X}
    \toprule
    \headerrow
    \thead{Estimador} & \thead{Fórmula} \\
    \midrule
    Media
      & $\equ{\bar{x}} = \eqop{\frac{1}{\eqk{n}}}
          \eqop{\sum}_{i=1}^{\eqk{n}} \eqk{x_i}$ \\
    Varianza
      & $\equ{s^2} = \eqop{\frac{1}{\eqk{n}-1}}
          \eqop{\sum}_{i=1}^{\eqk{n}} (\eqk{x_i} - \equ{\bar{x}})^2$ \\
    Error estándar
      & $\equ{\text{SE}} = \equ{s}/\sqrt{\eqk{n}}$ \\
    \bottomrule
  \end{tabularx}
\end{boxtab}
El estimador $\equ{s^2}$ es insesgado: $\eqop{\mathbb{E}}[\equ{s^2}] = \eqk{\sigma^2}$.
\end{definicion}

%───────────────────────────────────────────────────────────────────────────────
\section{Figuras y gráficas}
\label{sec:gal:figuras}
%───────────────────────────────────────────────────────────────────────────────

% CRÍTICO: nunca Ctp* en pgfplots — usar plot-series-N, plot-axis, etc.
% Estilos: catppuccin-clean, catppuccin-filled
% Cycle list: catppuccin (8 series, plot-series-1..8)

\begin{figure}[htbp]
  \centering
  \begin{tikzpicture}
    \begin{axis}[
      catppuccin-clean,
      width=0.85\textwidth, height=6cm,
      xlabel={Época},
      ylabel={Pérdida $\eqfn{\mathcal{L}}(\equ{\theta})$},
      title={Convergencia de optimizadores~\cite{goodfellow2016}},
      ymode=log, xmin=0, xmax=100,
      legend pos=north east, no markers,
    ]
      \addplot+[thick, domain=0:100, samples=80]
        {2*exp(-0.02*x)+0.10}; \addlegendentry{SGD}
      \addplot+[thick, domain=0:100, samples=80]
        {2*exp(-0.04*x)+0.05}; \addlegendentry{Momentum}
      \addplot+[thick, domain=0:100, samples=80]
        {2*exp(-0.07*x)+0.02}; \addlegendentry{Adam}
      \addplot+[thick, domain=0:100, samples=80]
        {2*exp(-0.05*x)+0.03}; \addlegendentry{RMSprop}
    \end{axis}
  \end{tikzpicture}
  \caption{Pérdida $\eqfn{\mathcal{L}}(\equ{\theta})$ por épocas.
           Colores del cycle list \texttt{catppuccin} (aliases
           \texttt{plot-series-1..4}, sin \texttt{Ctp*}).}
  \label{fig:gal:optimizers}
\end{figure}

\begin{figure}[htbp]
  \centering
  \begin{tikzpicture}
    \begin{loglogaxis}[
      catppuccin-clean,
      width=0.85\textwidth, height=6cm,
      xlabel={Paso de malla $\eqk{h}$},
      ylabel={Error $\|\equ{u}_h - \eqk{u}\|_2$},
      title={Convergencia FD centradas~\cite{leveque2007}},
      legend pos=north west,
      xtick={0.0625,0.125,0.25,0.5},
      xticklabels={$1/16$,$1/8$,$1/4$,$1/2$},
    ]
      \addplot[plot-series-1, mark=*] coordinates {
        (0.5,1.2e-2)(0.25,3.1e-3)(0.125,7.8e-4)(0.0625,1.9e-4)};
        \addlegendentry{FD centradas}
      \addplot[plot-series-5, dashed, no markers]
        coordinates {(0.5,1.2e-2)(0.0625,1.9e-4)};
        \addlegendentry{$\mathcal{O}(\eqk{h}^2)$}
    \end{loglogaxis}
  \end{tikzpicture}
  \caption{Convergencia $\mathcal{O}(\eqk{h}^2)$ del esquema de FD centradas
           para la ecuación de Poisson 1D.}
  \label{fig:gal:fd}
\end{figure}

\begin{figure}[htbp]
  \centering
  \begin{tikzpicture}
    \begin{axis}[
      catppuccin-clean,
      width=0.85\textwidth, height=5.5cm,
      xlabel={$\equ{x}$},
      ylabel={$\equ{\psi}_n(\equ{x})$},
      title={Eigenfunciones del oscilador armónico~\cite{griffiths2018}},
      domain=-4:4, samples=150, ymin=-1.2, ymax=1.5,
      no markers, legend pos=south east,
    ]
      \addplot+[thick] {0.7511*exp(-x^2/2)};
        \addlegendentry{$\equ{\psi_0}$}
      \addplot+[thick] {0.5311*2*x*exp(-x^2/2)};
        \addlegendentry{$\equ{\psi_1}$}
      \addplot+[thick] {0.2659*(4*x^2-2)*exp(-x^2/2)};
        \addlegendentry{$\equ{\psi_2}$}
      \addplot+[thick] {0.1083*(8*x^3-12*x)*exp(-x^2/2)};
        \addlegendentry{$\equ{\psi_3}$}
    \end{axis}
  \end{tikzpicture}
  \caption{Primeras eigenfunciones $\equ{\psi}_n = H_n e^{-\equ{x}^2/2}$
           del oscilador armónico cuántico.}
  \label{fig:gal:harmonic}
\end{figure}

%───────────────────────────────────────────────────────────────────────────────
\section{Diagramas TikZ}
\label{sec:gal:diagramas}
%───────────────────────────────────────────────────────────────────────────────

% CRÍTICO: nunca Ctp* en TikZ — usar tikz-*, env-*, fb-*, eq-* según el rol.
% tikz-initial (Blue): estado/nodo inicial
% tikz-accept (Green): estado aceptador / nodo destino
% tikz-node-border / background / maintext: nodo genérico
% eq-known (Mauve): pesos de aristas (datos conocidos del problema)
% env-phys (Sapphire): caja de función / transformación

\subsection{Autómata finito determinista}

\begin{figure}[htbp]
  \centering
  \begin{tikzpicture}[
    shorten >=1pt, node distance=2.5cm, on grid, auto, thick,
    state/.style={circle, draw=tikz-node-border, fill=background,
                  text=maintext, minimum size=1cm},
    every edge/.style={draw=tikz-edge, ->, >=latex},
  ]
    \node[state, initial, initial text={}, draw=tikz-initial]
      (A) {$\eqk{A}$};
    \node[state, right=of A] (B) {$\eqk{B}$};
    \node[state, right=of B] (C) {$\eqk{C}$};
    \node[state, accepting, draw=tikz-accept, right=of C]
      (D) {$\eqk{D}$};
    \path[->]
      (A) edge[loop above] node {$\eqk{a}$} (A)
          edge              node {$\eqk{b}$} (B)
      (B) edge[loop above] node {$\eqk{b}$} (B)
          edge              node {$\eqk{a}$} (C)
      (C) edge              node {$\eqk{b}$} (D)
          edge[bend left]   node[swap] {$\eqk{a}$} (A)
      (D) edge[loop right]  node {$\eqk{a},\eqk{b}$} (D);
  \end{tikzpicture}
  \caption{DFA para $L = \{\eqvec{w} \in \eqdom{\{a,b\}^*} \mid
           \eqvec{w}$ contiene $\eqk{ab}\}$~\cite{sipser2012}.
           Estado inicial \texttt{tikz-initial} (azul);
           aceptador \texttt{tikz-accept} (verde).}
  \label{fig:gal:dfa}
\end{figure}

\subsection{Grafo dirigido ponderado}

\begin{figure}[htbp]
  \centering
  \begin{tikzpicture}[
    node distance=2.5cm,
    vertex/.style={circle, draw=tikz-node-border, fill=tikz-node-bg,
                   text=maintext, minimum size=0.9cm, thick},
    every edge/.style={draw=tikz-edge, ->, >=latex, thick},
    wlabel/.style={text=eq-known, font=\small},
  ]
    \node[vertex, draw=tikz-initial, fill=tikz-initial!15] (s) {$\eqk{s}$};
    \node[vertex, right=of s]  (u) {$\eqk{u}$};
    \node[vertex, right=of u]  (v) {$\eqk{v}$};
    \node[vertex, below=of u]  (x) {$\eqk{x}$};
    \node[vertex, draw=tikz-accept, fill=tikz-accept!15,
          right=of x]          (t) {$\eqk{t}$};
    \path[->]
      (s) edge node[above, wlabel] {$\eqk{6}$} (u)
      (s) edge node[left,  wlabel] {$\eqk{7}$} (x)
      (u) edge node[above, wlabel] {$\eqk{5}$} (v)
      (u) edge node[right, wlabel] {$\eqk{-4}$}(t)
      (x) edge node[below, wlabel] {$\eqk{7}$} (t)
      (v) edge node[right, wlabel] {$\eqk{-2}$}(t);
  \end{tikzpicture}
  \caption{Grafo ponderado~\cite{cormen2022}: pesos en \texttt{eq-known} (malva)
           porque son datos del problema de camino mínimo.}
  \label{fig:gal:grafo}
\end{figure}

\subsection{Tabla hash con separate chaining}

\begin{figure}[htbp]
  \centering
  \begin{tikzpicture}[font=\small,
    cell/.style  = {draw=tikz-node-border, minimum width=1.3cm,
                    minimum height=0.65cm},
    occ/.style   = {cell, fill=tikz-accept,  text=background},
    empty/.style = {cell, fill=tikz-node-bg, text=maintext},
    fn/.style    = {draw=tikz-node-border, fill=env-phys, text=background,
                    minimum width=1.2cm, minimum height=3.2cm, align=center},
    key/.style   = {cell, fill=tikz-initial, text=background,
                    minimum width=2cm},
    arr/.style   = {draw=tikz-edge, ->, >=latex, thick},
  ]
    \node[key] (k1) at (-4.5, 0.9) {\texttt{k1}};
    \node[key] (k2) at (-4.5, 0.0) {\texttt{k2}};
    \node[key] (k3) at (-4.5,-0.9) {\texttt{k3}};
    \node[fn]  (hf) at (-2.4, 0)   {$\eqop{h}$};
    \foreach \i in {0,...,4} {
      \node[empty] (b\i) at (0.5, {0.9 - \i*0.65}) {};
      \node[left=0.1cm of b\i, text=tikz-label, font=\scriptsize] {$\i$};
    }
    \node[occ] at (b1) {\texttt{k1}};
    \node[occ] at (b3) {\texttt{k2}};
    \node[occ] (c1) at (2.1, {0.9 - 1*0.65}) {\texttt{k3}};
    \draw[arr] (b1) -- (c1);
    \draw[arr] (k1) -- (hf);
    \draw[arr] (k2) -- (hf);
    \draw[arr] (k3) -- (hf);
    \draw[arr] (hf.east) to[out=0,in=180] (b1.west);
    \draw[arr] (hf.east) to[out=0,in=180] (b3.west);
    \node[above=0.3cm of k1, text=maintext, font=\small] {\textbf{Claves}};
    \node[above=0.3cm of hf, text=maintext, font=\small] {\textbf{Hash}};
    \node[above=0.3cm of b0, text=maintext, font=\small] {\textbf{Buckets}};
  \end{tikzpicture}
  \caption{Separate chaining~\cite{cormen2022}: función hash $\eqop{h}$
           (\texttt{env-phys}, sapphire) mapea claves (\texttt{tikz-initial},
           azul) a buckets (verde=ocupado, gris=libre).}
  \label{fig:gal:hashing}
\end{figure}

\subsection{Diagrama conmutativo}

\begin{figure}[htbp]
  \centering
  \begin{tikzcd}[
    every arrow/.append style={draw=tikz-edge, thick},
  ]
    \eqdom{G}
      \arrow[r, "\eqop{\varphi}"]
      \arrow[d, "\eqop{\pi}"', two heads]
    & \eqdom{H} \\
    \eqdom{G}\,\eqop{/}\,\eqdom{\ker\varphi}
      \arrow[ur, dashed, "\eqop{\tilde{\varphi}}"', "\sim"']
  \end{tikzcd}
  \caption{Primer teorema de isomorfismo~\cite{mac1998}:
           $\eqdom{G}/\eqdom{\ker\varphi} \cong \eqop{\mathrm{im}}(\eqop{\varphi})$.}
  \label{fig:gal:isomorfismo}
\end{figure}

%───────────────────────────────────────────────────────────────────────────────
\section{Elementos auxiliares}
\label{sec:gal:auxiliares}
%───────────────────────────────────────────────────────────────────────────────

\dificultad{3}{4}

\begin{ejercicio}
Demuestre que la pérdida MSE para regresión lineal conduce a las ecuaciones
normales $\eqten{X}^\top\eqten{X}\equ{\hat{\beta}} = \eqten{X}^\top\eqk{y}$.
\end{ejercicio}

\begin{ejercicio}
Sea $\eqfn{f}(\equ{x}) = e^{-\equ{x}^2}$. Calcule la transformada de
Fourier $\eqop{\mathcal{F}}[\eqfn{f}]$ y verifique que es un punto fijo
de $\eqop{\mathcal{F}}$ (salvo escala).
\end{ejercicio}

\begin{notasbib}
La teoría de espacios de Hilbert y el análisis funcional riguroso siguen
a Reed y Simon~\cite{reed1980}. Los métodos de score matching y su variante
denoising~\cite{hyvarinen2005,vincent2011} son fundamentales en los
modelos de difusión modernos. La formulación variacional de EDPs y los
métodos de diferencias finitas se tratan en~\cite{leveque2007}.
\end{notasbib}

%───────────────────────────────────────────────────────────────────────────────
\printbibliography[heading=subbibliography]
